% ---------------------------------------------------------------------------
% Section 1: Introduction. Every number is a macro from tables/numbers.tex.
% ---------------------------------------------------------------------------

\section{Introduction}
\label{sec:intro}

When a recommender sits between the two sides of a matching market, it decides
which candidates anyone sees. An agent on nobody's shortlist is not rejected; it
is never considered, and no preference data record the difference. The usual way
to measure what the ranker adds is to compare exposure under the learned ranker
with exposure under an oracle that ranks by true utility, on the same population
and seed. The oracle's own concentration comes from preferences, since agents
that many others value highly appear on many lists; the difference between the
two is concentration the ranker adds.

In the market simulated here that difference changes with the horizon at which it
is read: large at first, small after a hundred rounds, and inside a margin of 0.002 Gini after three
hundred. A comparison at one horizon cannot tell a lasting bias
from a cost of the ranker's early training, and the two call for different
readings of what the ranker does. This paper measures the whole trajectory, how
it ends, and what sets its length.

\paragraph{The convergence objection.} The feedback-loop literature leads one to
expect the opposite. A system that learns from data its own choices generate can
lock in \citep{ensign2018,chaney2018,mansoury2020}, and a ranker that shows only
the agents it already rates highly might never learn about the others. The answer
in this model is a property of its design, stated here rather than discovered
later: the ranker's training data are not generated by its own shortlists. Each
round's log holds every agent's true top admissible candidates and uniform
negatives, so the ranker sees informative data about everyone, whether or not it showed them.
Once the gate stops binding, those positives are the oracle's own top-$k$ lists, so
the ranker is trained to imitate the benchmark it is scored against, and the gap
measures how fast that imitation converges. Its shortlists reach that log through one channel only: they decide who receives
proposals, which moves acceptance thresholds, which changes admissible sets.
\S\ref{sec:what-ends} measures how much that channel matters, and
\S\ref{sec:limits} says what changes in a market whose recommendations are its
only data.

\paragraph{Contributions.}
\begin{enumerate}
\item \textbf{A transient with a measured length.} On 40 seeds the gap peaks at
  $\peakGap$ near round \peakRound{}, and every seed enters and stays inside a
  pre-registered equivalence margin of $0.002$ between rounds \entryMinAll{} and
  \entryMaxAll{}. Over 1000 rounds an exponential decay predicts held-out rounds
  better than a power law or an exponential to a constant, with a half-life of
  \halfLife{} rounds and a residual of $\longTail$, inside the margin and above
  zero (\S\ref{sec:transient}, \S\ref{sec:floor}).
\item \textbf{What ends it.} At round 0 the untrained ranker's scores are almost
  all a per-candidate term, so every agent on a side sees nearly the same twenty
  candidates and the exposure Gini is the closed form $1 - 2k/N$. The gap closes
  as that common list breaks into each agent's own top twenty. Across training doses on fresh
  logs its length is set by optimizer steps, not rounds (elasticity
  $\doseElasticity$); a ranker trained on another market's data closes it slightly
  faster; and pretraining removes about four fifths of the decay on the registered
  measure and less on the other three (\S\ref{sec:what-ends}).
\item \textbf{An open market.} When a share of agents is replaced every round, the
  gap had not closed by round 600: over rounds 500 to 599 it is $\openGapLow$,
  $\openGapMid$ and $\openGapHigh$ at 0.5, 1 and 2 percent per round, still falling
  at the lowest rate and rising at the other two. At the two higher rates
  newcomers are under-exposed and long-tenured agents over-exposed relative to the
  oracle (\S\ref{sec:arrivals}).
\item \textbf{An institutional comparison that depends on the horizon.} Read at
  five rounds, no setting of a small complete-information subgraph beats the
  oracle-ranked market; read after the transient, a range of subgraph sizes at the
  lowest proposal cost beats it, and beats the learned-ranker market too
  (\S\ref{sec:institutional}).
\item \textbf{Sensitivity of single trajectories.} A $10^{-7}$ perturbation to the
  initial weights grows to the size of the between-seed spread, so a reported
  number estimates a property of the seed distribution, not of one run
  (\S\ref{sec:design}).
\end{enumerate}

\paragraph{What the model is about.} It is a two-sided matching market with a
learned ranker. One side's utility includes a bonus for a partner whose status
exceeds its own by a margin $\delta$. The exposure results do not depend on it:
across $\delta \in \{0, 0.5, 1\}$ the peak gap moves
by $\deltaPeakSpread$ and the median entry round by \deltaEntrySpread{} rounds,
against $\shockPeakSpread$ and \shockEntrySpread{} rounds across the chemistry
axis (\S\ref{sec:where}). The status bonus is the one asymmetry between the sides, so the welfare
comparison of \S\ref{sec:institutional} may depend on it; that dependence was not
tested.
Every parameter is a modeling choice and nothing is calibrated to data
(\S\ref{sec:limits}).
